\subsection{Тонкая линза}

%%%%%%%%%%%%%%%%%%%%%%%%%%%%%%%%%%%%%%%%%%%%%%%%%%%%
%%%%%%%%%%%%%%%%%%%%%%%%%%%%%%%%%%%%%%%%%%%%%%%%%%%%
 \z{Построить изображение точечного источника в тонкой линзе для случаев, изображённых на рисунке.
 		\begin{center}%
 			\includegraphics[width=1\textwidth]{4.5.1}%
		\end{center}%
		}
%
{}

%%%%%%%%%%%%%%%%%%%%%%%%%%%%%%%%%%%%%%%%%%%%%%%%%%%%
 \z{На рисунке показан ход светового луча 1 до и после линзы. Найти построением положение фокусов линзы и ход светового луча 2.
  		\begin{center}%
 			\includegraphics[width=1\textwidth]{4.5.2}%
		\end{center}%
 }
%
{}

%%%%%%%%%%%%%%%%%%%%%%%%%%%%%%%%%%%%%%%%%%%%%%%%%%%%
 \z{Светящаяся точка находится на главной оптической оси линзы с оптической силой $D = -2.5\,\text{дптр}$. Расстояние от линзы до её изображения $f = 30\,\text{см}$. На каком расстоянии от линзы находится точка?}
%
{$d = 120\,\text{см}$}

%%%%%%%%%%%%%%%%%%%%%%%%%%%%%%%%%%%%%%%%%%%%%%%%%%%%
 \z{Предмет помещают на главной оптической оси собирающей линзы на расстоянии $d = 20\,\text{см}$ от неё и получают действительное изображение предмета на расстоянии $f = 4F$, где $F$ — фокусное расстояние линзы. Определить фокусное расстояние.}
%
{$F = 15\,\text{см}$}

%%%%%%%%%%%%%%%%%%%%%%%%%%%%%%%%%%%%%%%%%%%%%%%%%%%%
 \z{Предмет помещают на главной оптической оси рассеивающей линзы на расстоянии $d = 1.5F$, где $F$ — фокусное расстояние линзы. Изображение предмета при этом получается на расстоянии $f = 20\,\text{см}$ от линзы. Определить фокусное расстояние линзы.}
%
{$F = -33.3\,\text{см}$}

%%%%%%%%%%%%%%%%%%%%%%%%%%%%%%%%%%%%%%%%%%%%%%%%%%%%
 \z{Какое минимальное расстояние между предметом и его действительным изображением для линзы с фокусным расстоянием $F$?}
%
{$4F$}

%%%%%%%%%%%%%%%%%%%%%%%%%%%%%%%%%%%%%%%%%%%%%%%%%%%%
 \z{Расстояние между точечными источниками света и экраном $L = 3\,\text{м}$. Линза, помещённая между ними, даёт чёткое изображение при двух положениях, расстояние между которыми $l = 1\,\text{м}$. Найти фокусное расстояние линзы.}
%
{$F = 67\,\text{см}$}
